\scanpage{65}

\mspage{6}

(i) $M$ est \uline{plat} sur $A$

(ii) $M$ est \uline{formellement projectif} sur $A$

\uline{démonstration} : a) réduction à $A$ artinien local

\[
\begin{aligned}
&\text{b)}\quad\mathrm{(ii)}\Longleftrightarrow\quad\forall N\text{ sur }A,\quad\varinjlim\operatorname{Ext}^1_A(M/n^kM,N)=0\Longleftrightarrow\quad\text{itou pour }N\text{ un}\\
&k\text{-vectoriel (dévissage)}\Longleftrightarrow\quad\varinjlim\bigl(\operatorname{Ext}^1_A(M/n^kM,k)\bigr)^I=0\quad\forall I\\
&\Longleftrightarrow\quad\text{``}\varinjlim\text{''}\operatorname{Ext}^1_A(M/n^kM,k)=0\ \text{comme Ind-objet}\Longleftrightarrow\quad\text{``}\varprojlim\text{''}\operatorname{Tor}^1_A(M/n^kM,k)=0
\end{aligned}
\]

comme pro-objet, le précédent étant dual de celui-ci. On calcule ce $\operatorname{Tor}$ par une résolution plate de $k$, et par Artin-Rees, le pro-objet précédent n'est autre que ``$\varprojlim$'' $\operatorname{Tor}^1_A(M,k)/n^k\operatorname{Tor}^1_A(\ldots)$ dont la nullité signifie $\operatorname{Tor}^1(M,k)=0$. i e que $M$ est plat sur $A$.

$0_{\overline{\underline{\mathrm{IV}}}}$ 19.4.6\qquad (changement de base et lissité formelle) - Il n'est pas nécessaire de supposer $B$ formellement projectif dans l'hypothèse 2°, en faisant sur la classe de l'extension le jeu fait ici sur le cocycle de Hochschild - De même, dans 1°, l'hypothèse de ``formellement projectif'' peut se remplacer par une hypothèse de finitude ($B/J_{\unsure{\lambda|s}}$ fini sur $A/J_\lambda$ qui est noethérien p.e.)

\margin{?}

$0_{\overline{\underline{\mathrm{IV}}}}$ 19.4.7.\qquad La démonstration se simplifierait si tu avais fait remarquer auparavant que $\operatorname{\unsure{Exalcotop}}_A(B,L)$ est foncteur semi-exact en $L$.

$0_{\overline{\underline{\mathrm{IV}}}}$ 19.8.3\quad
$\left\{\begin{array}{ll}\text{pg 110 ligne}&-8\\&-6\\\text{pg 111 ligne}&1\end{array}\right.$ : lire homomorphisme local \quad et non \quad homomorphisme

$0_{\overline{\underline{\mathrm{IV}}}}$ 20.4.6\qquad Le signe choisi pour $dx$ est l'opposé du signe universellement choisi en analyse (où $\Delta f(x;x')=f(x')-f(x)$, $g.\Delta f(x;x')=g(x)\,\Delta f(x;x')$), compte tenu de ce que $B\otimes B$ et $P^1_{B/A}$ sont des $B$-\deleted{modules}\inserted{algèbres} par $b\longrightarrow b\otimes1$.

% slanted
\margin{faire errata ; \unsure{remplacer} $p_1-p_2$ par $p_2-p_{\unsure{2|1}}$ ! attention aux \unsure{passages|pass.} subséquents !}

$0_{\overline{\underline{\mathrm{IV}}}}$ 20.4.13 (iv)\qquad L'hypothèse ``$A$ intégralement clos'' n'a rien à faire ici [cf 20.5.9 par ex.]

\margin{!}

enfin

$\overline{\underline{\mathrm{III}}}$ 4.3.6 pg 133 ligne 7 \quad la formule $K'=K\otimes_A\widehat A$ est fausse en général !

\margin{!!!}

En relisant cette lettre, je suis honteux de la démonstration tordue donnée pg 2 de ce que $\operatorname{Tor}$ et ``pro-objet complétion'' commutent. Il suffirait, par dévissage, de traiter directement le cas $N=A$ (et $M=B$), ié ``$\varprojlim$'' $\operatorname{Tor}^p_A(A/m^k,M)=0$ si $p>0$, mais j'en suis incapable. Tu trouveras ci-joint une lettre de Serre que je te renvoie

Bien à toi\\
P. Deligne\\
% archive sheet number: 64
